\documentclass{article}
\usepackage{PRIMEarxiv} % Core package for the PrimeAI Template
\usepackage{amsmath, amssymb, amsthm, bm, dcolumn} % Add amsthm here for the proof environment
\usepackage[numbers,sort&compress]{natbib} % Natbib for citations
\usepackage{graphicx} % For high-quality images
\usepackage{multicol} % Multi-column figures/tables
\usepackage{hyperref} % Hyperlinks
\usepackage{listings} % Code listings
\usepackage{authblk} % For structured affiliations
% Define theorem style (optional)
\theoremstyle{plain}
\newtheorem{theorem}{Theorem}
\renewcommand{\qedsymbol}{}

% Custom commands for primes and qubit encoding
\newcommand{\primeQubit}[1]{|p_{#1}\rangle}
\newcommand{\primeGate}[1]{U_{p_{#1}}}

\usepackage{wrapfig}
\usepackage[pscoord]{eso-pic}
\usepackage[fulladjust]{marginnote}
\reversemarginpar

% Typesetting improvements without footnote patching
\usepackage[protrusion=true, expansion=true, tracking=false]{microtype}
\microtypecontext{spacing=nonfrench}

% Line numbers
\usepackage[right]{lineno}

% Text layout - adjust as needed
\raggedright
\setlength{\parindent}{0.5cm}
\textwidth 5.25in 
\textheight 8.75in

% Set double spacing
\usepackage{setspace} 
\doublespacing

% Adjust width for specific content
\usepackage{changepage}

% Adjust caption style
\usepackage[aboveskip=1pt,labelfont=bf,labelsep=period,singlelinecheck=off]{caption}

% Remove brackets from references
\makeatletter
\renewcommand{\@biblabel}[1]{\quad#1.}
\makeatother

% Header, footer, and page numbers
\usepackage{lastpage,fancyhdr}
\pagestyle{fancy}
\fancyhf{}
\fancyhead[L]{Preprint - PrimeAI Enhanced Template}
\fancyfoot[C]{\scriptsize Multiplicity Theory © 2024 Ryan Van Gelder - Citizen Gardens \\ Licensed Under MIT and CC BY-NC-SA 4.0.}
\fancyfoot[R]{Page \thepage\ of \pageref{LastPage}}
\renewcommand{\footrule}{\hrule height 2pt \vspace{2mm}}

\begin{document}
\title{Multiplicity-Driven Temporal Quantum States: A Theoretical Framework and Applications}

\author{Ryan Van Gelder}
\affil{Citizen Gardens - The Foundation of Multiplicity}
\date{\today}
\maketitle

\begin{abstract}
This paper explores the theoretical underpinnings and practical applications of multiplicity-driven temporal quantum states. By introducing time-dependent multiplicity operators, we examine the possibility of simulating "time-split" quantum systems where particles simultaneously occupy states across different moments in time. Applications range from simulating time travel and retrocausal systems to advancing quantum computing and artificial general intelligence (AGI).
\end{abstract}

\section{Introduction}
Multiplicity theory emphasizes the interconnectedness of all elements in the universe. By extending its principles to time-dependent quantum systems, we propose a framework where particles can occupy multiple states across time, governed by time-evolving multiplicity equations. This builds upon foundational works in quantum mechanics and multiplicity theory \cite{vanGelder2024Multiplicity}. 

\section{Mathematical Framework}
The dynamics of temporal quantum states are modeled using the enhanced dynamic multiplicity equation:
\begin{align}
\frac{\partial \rho_k}{\partial t} &= \alpha_k(t) \rho_k + \beta_k(t) I_k + \gamma_k(t) \sum_j T_{kj} \rho_j \\
&+ \lambda(t) \Omega(\rho) + \zeta_k \sum_{m,n} C_{mn} \rho_m \rho_n + \xi_k(t),
\end{align}
where:
\begin{itemize}
    \item $\rho_k$ represents the state density at time $t$.
    \item $\alpha_k(t), \beta_k(t), \gamma_k(t)$ are time-dependent parameters.
    \item $\Omega(\rho)$ encodes geometric feedback dynamics.
    \item $\zeta_k \sum_{m,n} C_{mn} \rho_m \rho_n$ represents quantum entanglement terms.
    \item $\xi_k(t)$ introduces stochastic noise.
\end{itemize}
This equation allows for the exploration of time-reversed systems and temporal coherence through prime-based encoding and tensor networks \cite{galioto2024}.

\section{Simulating Time-Split Quantum Systems}
Temporal multiplicity offers a pathway to simulate retrocausal phenomena and time-reversed states. Using recursive feedback mechanisms, time-evolved states $\rho_k(t)$ can influence their past and future states, leading to potential applications in:
\begin{itemize}
    \item \textbf{Quantum Computing:} Enabling time-coherent memory systems.
    \item \textbf{Cosmology:} Modeling black hole dynamics and cosmic inflation.
    \item \textbf{Artificial Intelligence:} Enhancing neuromorphic architectures with time-aware adaptations.
\end{itemize}

\section{Applications}


\section{Challenges and Future Directions}
\subsection{Computational Complexity}
The non-linear and stochastic nature of the equations necessitates advanced computational models. Tensor networks and hybrid quantum systems can mitigate some challenges.

\subsection{Experimental Validation}
Developing experimental setups to validate temporal coherence and multiplicity remains a priority. Quantum processors could simulate time-split states in controlled environments.

\section{Conclusion}
Multiplicity-driven temporal quantum states provide a theoretical and computational framework for addressing complex quantum phenomena and advancing interdisciplinary applications. Future research should focus on refining algorithms, experimental validation, and scaling applications to practical domains.


\nocite{*}
\bibliographystyle{plain}
\bibliography{references}

\end{document}